\section[Split Heegner determinants and paired logarithmic measures]{Split Heegner determinants and\\paired logarithmic measures}\label{sec:split-measures}

The split-pair anchor requires an integral comparison over an imaginary
quadratic field.  We first construct its strict determinant and paired
logarithmic measure.  Section~\ref{sec:split-local} compares them and proves
the following proposition.  The variables \(t,u\) and the auxiliary primes
introduced below are independent of those in the positive comparison.

\begin{proposition}[Reducible split-product comparison]
\label{split-measures:reducible-input}
Let \(E/\mathbb Q\) be a non-CM elliptic curve with
\(E(\mathbb Q)[2]\ne0\), and put \(N=\mathrm{cond}(E)\).
Let \(d\equiv7\pmod {16}\) be prime and suppose that every prime dividing
\(2N\) splits in \(K=\mathbb Q(\sqrt{-d})\).  If \(s_2(E/K)=1\), then
\(L(E/K,s)=L(E,s)L(E^{-d},s)\) has a simple zero at \(1\), and
\[
                              X(E)+X(E^{-d})=0 .
\]
If the product is already known to have a simple zero at \(1\), the same
valuation identity holds without separately assuming the Selmer corank
condition.
\end{proposition}

The superscript \(-d\) denotes a quadratic twist.  In the corank-one
assertion, we do not assume the existence of a nontorsion point: this will
follow from the determinant comparison at the tame center.

\subsection{The strict determinant and its residual order}

Until explicitly enlarged below, the setting is the geometric one in
Proposition~\ref{split-measures:reducible-input}: \(E\) is non-CM with
rational two-torsion, \(d\) is prime and congruent to \(7\pmod {16}\),
and every prime of \(2N\) splits in \(K\).  The constructions and
lemmas in this section impose no Selmer-corank hypothesis.

We fix embeddings for CM and 2-adic evaluations and a place \(w\) over 2 (the connected prime in our CM convention), writing \(\bar w\) for the conjugate. Let \(\mathcal O_j=\mathbb Z+j\mathcal O_K\) with ring class field \(H_j\), \(h=h(K)\), and \(\Gamma\simeq\mathbb Z_2\) the maximal free pro-2 quotient of \(\varprojlim_m\operatorname{Pic}(\mathcal O_{2^m})\). We use class groups as Galois labels by ring class reciprocity. Take \(\Psi_t\) the tautological character on \(\Gamma\) with generator value \(1+t\), pulled back to a Galois character over \(K\). 

Here \(h\) is odd by quadratic genus theory, and \(\mathcal O_K^\times=\{\pm1\}\). The order class group exact sequence uses, at 2, the ratio between the unit groups at the two primes modulo the order conductor. This gives the description of \(\Gamma\) and surjective inertia at either dyadic place. 

At a split odd prime \(q\), the Frobenius exponents at its two places in \(\Gamma\) are opposite and nonzero: taking the \(h\)-th power of a place ideal uses the ratio of a principal generator and its conjugate, which is not a root of unity. Conjugation in all our ring class character groups acts by inversion.

Use \(T=T_2E\), with Tate duality given by the Weil pairing.  The strict
complex is full at all allowed places except \(w\):
\[
 C(t)=\mathrm{fib}\big(R\Gamma(G_{K,S},T[[t]](\Psi_t))
                        \longrightarrow R\Gamma(K_w,T[[t]](\Psi_t))\big),
 \qquad S=\{v:v\mid 2N\}.
\]
All complexes and limiting operations use the finite models of
Section~\ref{sec:foundations}.  The strict complex has a square free model in
degrees \(1,2\); the same assertion will hold for its two-variable analogue.
Indeed the residual representation has two trivial constituents.  For each
constituent, restriction to \(w\) kills global degree-zero constants, and
the complementary restriction detects dual constants, giving no degree-three
cohomology.  The full local complex at \(w\) subtracts exactly the global
Euler-characteristic contribution.  The assertions for \(T/2T\) follow by
its filtration and lift to the free model.

Write \(L_{\rm alg}(C)\) for the determinant of the model differential, defined
up to a coefficient-ring unit.  At a valuation where it is nonzero, our
inverse-determinant convention gives
\[
 d(C,1)=-v\bigl(L_{\rm alg}(C)\bigr).
\]

\begin{lemma}[Scalar residual order]
\label{split-measures:residual-order}
For an odd prime \(q\mid N\), let \(b_q\in\mathbb Z_2\setminus\{0\}\) be
the \(\Gamma\)-Frobenius exponent at one of the places over \(q\), and put
\(g_q=2^{v_2(b_q)}\).  Then
\[
 \operatorname{ord}_t\overline L_{\rm alg}(C(t))
   =4\sum_{\substack{q\mid N\\q\ne2}}g_q.
 \tag{A1}\label{eq:A1}
\]
\end{lemma}
\begin{proof}
Let \(w^h=(\alpha)\) as ideals, \(\alpha=(a+b\sqrt{-d})/2\), \(a^2+db^2=2^{h+2}\). Both \(a,b\) are odd: otherwise \(\alpha\) would have the same reduction at both primes over 2, impossible. Each odd prime of \(b\) has \((2/q)=1\) (here \(q\) denotes that prime), by the norm equation and \(h\) odd; hence \(b\equiv\pm1\bmod8\). For \(h\ge3\), \(a^2\equiv9\bmod16\), and the unit \(\bar\alpha\) at \(w\) is congruent to \(a\) modulo 8. For \(h=1\) we have \(d=7\), and the unit root of \(X^2\pm X+2\) is again \(\pm3\bmod8\). 

Thus \(-1,2,\bar\alpha\) give a basis of local squareclasses there. They are also a basis of global squareclasses with only dyadic support (odd \(h\)). With trivial scalar \(\mathbb F_2\) and no odd places yet allowed, the strict problem is acyclic: restriction to \(w\) is an isomorphism in degree one by that basis and in degree zero on constants, so use the above amplitude and Euler characteristic.

On either individual residual constituent with scalar twist \(\overline\Psi_t\), now adding the two places over each odd \(q\mid N\) gives singular blocks by the unramified inflation comparison. Each has determinant over \(\mathbb F_2[[t]]\) of order \(g_q=2^{v_2(b_q)}\) in \(t\), \(b_q\in\mathbb Z_2\) the exponent at one place over \(q\). Indeed Frobenius-minus-one on the scalar inertia \(H^1\) coefficient is \((1+t)^{\pm b_q}-1\) (the cardinality twist reduces to one). Extension of the two constituents is only needed at full support. There are two places over each \(q\) and two residual constituents, so
summing these orders gives \eqref{eq:A1}.  Determinant multiplicativity for
the residual filtration allows the two constituents to form a nonsplit
extension.
\end{proof}

Retain the generator \(\alpha\) of \(w^h\) chosen in the proof.

\subsection{A transverse tame parameter}

\begin{lemma}[Auxiliary primes and Frobenius exponents]
\label{split-measures:auxiliary-primes}
There is a sequence of split good primes \(r_i\), leaving every fixed finite
set, with principal places \(R_i=(\pi_i)\), such that:
\begin{enumerate}[label=\textup{(\roman*)}]
\item \(\pi_i\) tends to \(1\) at both dyadic places, and its residue modulo
\(\sqrt{-d}\) is a nonzero nonsquare;
\item \(\sqrt2\) exists and is a nonsquare modulo \(R_i\), and
\((q/r_i)=-1\) for every odd \(q\mid N\);
\item the Tate Frobenius at \(R_i\) tends to a matrix \(g\) with
\(\det g=1\) and \(a_*:=\operatorname{tr}(g)\ne2\).
\end{enumerate}
The Sylow-2 ring-class character at conductor \(r_i\) has odd Frobenius
exponents at the odd places over \(N\), and exponents of valuation exactly
one at \(w,\bar w\).
\end{lemma}
\begin{proof}
For existence, prescribe a compatible ray class action over \(K\) of moduli \(2^m\sqrt{-d}\), using principal ideals with generators having residues 1 at the 2's and a fixed nonsquare at the ramified prime. This fixes 2-power roots of unity by norm. \(K(i,\sqrt[4]{2})/K\) is dihedral of order 8 (disjoint quadratic data); its maximal abelian subfield is 2-cyclotomic in the sense of being contained in \(K(\mu_{2^\infty})\). Thus its nontrivial commutator element can be prescribed compatibly with the ray condition. The \(\sqrt q\)'s are independently switchable by their new disjoint ramification. 

The full constraint field has finite commutator image over \(K\), so the subgroup fixing it has Tate image containing a finite-index part of the closed commutator of the image over \(K\). This includes an open subgroup in determinant one by the non-CM open-image theorem~\cite{Serre72}. The coset therefore permits trace different from 2. After fixing such an element, Chebotarev to increasing precisions (degree-one places with chosen lifts) gives the sequence; in particular \(r_i\to1\) 2-adically.

The Sylow-2 quotient of \(\operatorname{Pic}(\mathcal O_{r_i})\) is cyclic of order \(2^{v_2(r_i-1)}\) by the unit ratio at \(r_i\) and odd class number. Use its group-ring Galois character \(\Psi_{u,i}\) with generator value \(1+u\). Its quadratic quotient is \(K(\sqrt{r_i})\) (pull back the discriminant character by norm, trivial on the defining principal order relations). 

At the places over odd \(q\mid N\), Frobenius exponents in this variable are odd: for \(\mathfrak q^h=(x)\), the ratio \(x/\bar x\bmod R\) has symbol \((q/r)=-1\). 

At \(w,\bar w\) the exponents have valuation exactly one. Indeed \(\alpha\bmod R\) is square by Hilbert reciprocity on \((\alpha,\pi_i)\), primary \(\pi_i\) at both 2's and units at other odd places. Thus \(\alpha/\bar\alpha\) is a square with nonsquare root \(\alpha/(\sqrt2)^h\bmod R\); both roots have the same symbol. Inertia at either place over \(r\) surjects onto the cyclic quotient, while \(\Psi_t(\mathrm{Fr}_R)\to1\) since \(\pi_i\) is increasingly primary.

\end{proof}

Fix the sequence from Lemma~\ref{split-measures:auxiliary-primes} and
retain its generator choices and group-ring characters \(\Psi_{u,i}\).
Put \(n_i=v_2(r_i-1)\) and
\[
 \Lambda_i=\mathbb Z_2[[t]][u]/\big((1+u)^{2^{n_i}}-1\big).
\]
At stage \(i\), let \(C_i\) be the strict complex with coefficients
\(T\otimes\Lambda_i\), scalar action \(\Psi_t\Psi_{u,i}\), and the
two places over \(r_i\) additionally allowed in full.  It is still strict
at \(w\).  Since \(n_i\to\infty\), the displayed relation disappears
at every fixed Artin precision.  The finite-model construction on the
fixed ultrafilter therefore gives a matrix limit over
\(\mathbb Z_2[[t,u]]\), which we denote by \(C_r(t,u)\), with determinant
\(L_{\rm alg,r}\).  Thus the subscript \(r\) on a limiting object denotes
the chosen auxiliary sequence; in a stage calculation we write
\(r=r_i\) and \(R=R_i\).

\begin{lemma}[Transverse regularity]
\label{split-measures:transverse}
The determinants satisfy
\[
                   L_{\rm alg,r}(t,0)
                   =(2-a_*)^2 L_{\rm alg}(C(t))\cdot{\rm unit},       \tag{A2}\label{eq:A2}
\]
and
\[
 L_{\rm alg,r}\big|_{(2,t)=0}\ne0.
 \tag{A3}\label{eq:A3}
\]
\end{lemma}
\begin{proof}
At \(u=0\), unramified inflation supplies two singular local differentials
with Tate matrices \(g-1\), up to conjugation and determinant units.  The
cardinality twist tends to one, and \(\det(g-1)=2-a_*\).  This proves
\eqref{eq:A2}.

For \eqref{eq:A3}, set \(2=t=0\) and consider each trivial residual
constituent with its \(u\)-scalar twist over
\(D_0=\mathbb F_2[[u]]\). Discard the odd \(q\mid N\) for generic acyclicity on the constituent, by the singular inflation comparison using the nontrivial unramified scalar there.

In the remaining strict problem, at \(u=0\) degree one injects into unrestricted cohomology and has basis \(\pi_i,\bar\pi_i\) (radical classes, primary at \(w\), using the preceding dyadic basis). This notation identifies such classes by the ultrafilter. 

Suppose the two-term differential over \(D_0\) has a nontrivial free kernel; its saturation gives \(c\in H^1\) reducing nontrivially. Lifting modulo \(u^2\) forces the global cup of that reduction with \(\chi=[r_i]\), the linear scalar character coefficient, to vanish. 

At \(R\) both basis elements pair nontrivially with \(\chi\). Indeed \(r\equiv1\bmod8\), and \(\bar\pi_i\bmod R\) has symbol \((2B/r)=-1\), \(2B=\operatorname{Tr}(\pi_i)\). Here \(B\) is an integer, \(B\equiv1\bmod4\) by primarity, \(4r=(2B)^2+d\gamma^2\); at every odd prime \(q\mid B\) we have \((r/q)=(d/q)\) with both units (norm and the ramified-prime condition). Quadratic reciprocity also for possible \(B<0\) gives \((B/r)=(B/d)=-1\). Since the local \((r,r)\) Hilbert symbol is trivial and \(r=\pi_i\bar\pi_i\), this proves the assertion. Hence \(c\) reduces in global cohomology to \(\chi\).

In unrestricted global cohomology put
\[
 D_\chi=\frac{\Psi_u-1}{u}.
\]
This is a twisted cocycle: at a finite stage it is defined by retaining one
additional coefficient precision.  It reduces to \(\chi\).  The coefficient
exact sequence for multiplication by \(u\) therefore gives an integral
class \(z\) such that
\[
 [c]-[D_\chi]=u[z].
\]
Only existence of this quotient is needed; unrestricted \(H^1\) need not be
\(u\)-torsion-free.

At \(w\), \(D_\chi+uz\) is a boundary, and \(\Psi_u=1+u^2{\rm ur}\bmod u^3\) with \({\rm ur}\) the nontrivial unramified quadratic additive character. Thus \(\bar z_w={\rm ur}\) (boundaries are divisible by \(u^2\)). 

At \(\bar w\), lifting \(\bar z\) modulo \(u^3\) with the same leading coefficient forces \(\bar z_{\bar w}\cup {\rm ur}=0\), hence even valuation as a radical. 

These linear-term obstruction and boundary calculations hold on actual cochains to fixed Artin precisions by the finite models; fixed dyadic local actions here stabilize modulo each precision in the ultrafilter sense. The available global radical classes are generated by
\[
 -1,\quad 2,\quad\bar\alpha,\quad\pi_i,\quad\bar\pi_i.
\]
The class \(\bar z\) has even valuations at both dyadic places: at \(w\)
this follows because its image is unramified, and at \(\bar w\) from the
cup obstruction.  These two valuation conditions remove the generators
\(2,\bar\alpha\), since \(h\) is odd.  Thus only
\(-1,\pi_i,\bar\pi_i\) remain.  The unramified class at \(w\) is the
radical class of \(5\),
whereas \(\pi_i,\bar\pi_i\) are local squares and the remaining image is
contained in the line generated by \(-1\).  This contradiction rules out a
free kernel on either scalar constituent, proving \eqref{eq:A3}.
\end{proof}

\subsection{Integral depleted primitives on the ordinary locus}

The next three subsections construct the measures that will be compared
with these determinants.  Their scope is broader than the residual
calculations above: the primitive, disk-measure, and interpolation
constructions apply to a primitive weight-two newform of trivial
character and an imaginary quadratic field with units \(\{\pm1\}\),
split at \(2\), at the primitive level, and at the odd support
\(S^\circ\) specified below.
The modular-abelian-quotient variant uses a fixed characteristic-zero
coefficient projection.  We return to the rational-two-torsion and
prime-discriminant hypotheses in the final trace and residual-order
calculation.

The use of depleted primitives to express CM sums in terms of Heegner
logarithms has its antecedents in
Bertolini--Darmon--Prasanna~\cite[Sections~3.8 and~5.3]{BertoliniDarmonPrasanna2013}.
We give the integral dyadic construction with the local and
parametrization factors used in our comparisons.

Use a modular parametrization \(\phi:X_0(N)\to E\), \(\phi(\infty)=0\), \(\phi^*\omega_E=c_E f\,dq_{\rm Tate}/q_{\rm Tate}\) for the normalized primitive form \(f\). We give the ordinary-locus calculation also in the variant for a primitive weight-two newform of trivial character: there use Abel--Jacobi from its actual level with base cusp, projected to the corresponding modular abelian quotient up to isogeny, and the logarithm functional pulling back to \(f\,dq_{\rm Tate}/q_{\rm Tate}\) at the chosen coefficient embedding. Denote this value function by \(G\) on CM disks, so \(G=c_E^{-1}\log_{\omega_E}\phi\) for \(E\). Cuspidal differences are torsion~\cite{ManinDrinfeld} and hence are killed by the logarithm. Coefficient extensions and comparisons using \(G\) itself here are in characteristic zero.

Let \(S^\circ\) consist of the odd rational primes of the chosen support (for the current input, the odd primes dividing \(N\)); for the measure construction it can include additional split primes. Write
\[
 \begin{gathered}
 P_\ell(Z)=1-\frac{a_\ell}{\ell}Z+\frac{\epsilon_\ell}{\ell}Z^2,\\
 \epsilon_\ell=
 \begin{cases}
 1,&\ell\text{ does not divide the primitive level},\\
 0,&\ell\text{ divides the primitive level}.
 \end{cases}
 \end{gathered}
\]
where \(a_\ell=a_\ell(f)\). Fix split tame cyclic level orientations at the CM points, extending them to raised powers for quotient operators when needed. At 2 use the canonical connected cyclic level only. With \(V_\ell\) normalized by Tate expansion \(q_{\rm Tate}\mapsto q_{\rm Tate}^\ell\), form
\[
 \begin{split}
 f^{[2]}&=(1-a_2 V_2+2\epsilon_2 V_2^2)f,\\
 F_0&=d_{\rm mod}^{-1}f^{[2]}=\lim_{j\to\infty} d_{\rm mod}^{\,2^j-1} f^{[2]},\qquad
 F=\prod_{\ell\in S^\circ} P_\ell(V_\ell)\,F_0 .
 \end{split}
\]
Here \(d_{\rm mod}\) is ordinary modular differentiation; on a trivialized
Tate chart it is \(q_{\rm Tate}\partial_{q_{\rm Tate}}\).

\begin{lemma}[Integral depleted primitive]
\label{split-measures:integral-primitive}
The displayed limit defining \(F_0\) converges to an integral weight-zero
function on the ordinary locus, including when \(2\mid N\), and
\(d_{\rm mod}F_0=f^{[2]}\).  At a Tate cusp with cyclic \(\mu\)-level,
\[
 F_0=\sum_{\substack{m\ge1\\2\nmid m}}\frac{a_m}{m}q_{\rm Tate}^m,
 \qquad
 F=\sum_{\substack{m\ge1\\(m,\,2\prod_{\ell\in S^\circ}\ell)=1}}
       \frac{a_m}{m}q_{\rm Tate}^m.
\]
Congruences between such fully depleted primitives can be checked by their
Tate expansions on a common raised tame level, also for tuples of
coefficient components relative to a fixed coefficient lattice.
\end{lemma}
\begin{proof}
Work on the smooth ordinary formal curve of a fine full odd auxiliary level sufficiently large for the tame data, over \(\mathcal V=W(\overline{\mathbb F}_2)\) (and extend constants as needed). Use the canonical finite flat \(2\)-level from the multiplicative part, and the étale Igusa tower trivializing the multiplicative formal group and its differential. We use the ordinary deformation and tame elliptic moduli theory of~\cite{DeligneRapoport73,KatzMazur85}: the multiplicative part lifts over the ordinary deformations (its dual is étale); these trivialization covers extend étale on the ordinary compactified charts, with the multiplicative trivialization on Tate charts. 

The log Kodaira--Spencer differential given by the square of the trivialized differential gives \(d_{\rm mod}\) by its dual derivation on scalar functions; derivations lift along the étale levels and the completion. This uses the elliptic Kodaira--Spencer isomorphism on the fine prime-to-2 curve, with \(dq_{\rm Tate}/q_{\rm Tate}\) normalization and unit cusp widths there. It raises equivariant weights by two. \(V_2\) uses the quotient by the first connected subgroup with canonical choices on the result (for weight computations take target differential with pullback twice the source differential); \(V_\ell\) similarly uses the first subgroup in raised cyclic \(\ell\)-level with the stated Tate normalization.

Pullbacks of the algebraic characteristic-zero forms along the indicated ordinary maps have bounded denominator: these are rigid sections on the quasi-compact ordinary formal region, holomorphic also on cusp charts (the connected level and quotient on Tate curves give the usual Tate-level choices). One can first bound on a finite formal-affine covering in the Hodge lattice, before trivializing. 

Integrality, once some bound exists, can be checked by expansion at all trivialization lifts of a \(\mu\)-cyclic test cusp in each tame determinant component. Indeed the prime-to-2 full level special curve is smooth with geometrically connected determinant components. Each special Igusa component at any fixed finite height surjects onto its ordinary component, and formal expansions there detect vanishing of the special-fiber section. An integral section on the completed infinite tower uses just finite height modulo each precision. This gives the integrality test by division one uniformizer factor at a time. Choose the full-level cusps so that all the nested cyclic tame choices involved are of \(\mu\)-type. 

The weights of the inverse-derivative sequence tend 2-adically to zero, so expansion weight factors on all differential trivializations tend uniformly to 1. Its odd-exponent expansion converges uniformly modulo each scalar precision and the even terms vanish. The same expansion test on differences gives integral convergence on the tower; invariance of the limit descends it to weight zero (étale torsor descent levelwise modulo finite precisions). Also \(d_{\rm mod}F_0=f^{[2]}\). The Hecke recursions give the displayed depleted expansions.  The same
argument checks lattice congruences on a common raised tame level.  For a
tuple of coefficient components, first allow a common initial denominator
and then apply the argument in coordinates for the prescribed lattice.
\end{proof}

For the reducible input we also use the combination at \(1,V_2,V_2^2\) with coefficients \(1,-3,2\) of \(-1/24+\sum_{n\ge1}\sigma_1(n)q_{\rm Tate}^n\), the weight-two modular Eisenstein 2-depletion (the usual level-raised combination of the weight-two Eisenstein series). Denote its weight-zero antiderivative by \(F_{{\rm Eis},0}\). Since \(a_n\equiv\sigma_1(n)\bmod2\) off \(2N\) by residual Frobenius and Hecke recursions, \(F\bmod 2\) equals the fully depleted Eisenstein function, using \(1-(\ell+1)V_\ell/\ell+V_\ell^2/\ell\) at each \(\ell\in S^\circ\) on \(F_{{\rm Eis},0}\).

\subsection{Disk measures and exact logarithmic identities}

On an ordinary Serre--Tate disk choose the multiplicative coordinate \(z\)
with \(z=1\) at the canonical lift.  We use \(\mathcal F(z)\) for the disk
expansion of any of the weight-zero 2-depleted functions above; this
notation distinguishes the disk expansion from the global function.

\begin{lemma}[Support of the disk measure]
\label{split-measures:unit-support}
There is a bounded integral measure \(\mu\) on \(\mathbb Z_2\) with
\[
 \mathcal F(z)=\int_{\mathbb Z_2}z^x\,d\mu(x).
\]
It is supported on \(\mathbb Z_2^\times\).
\end{lemma}
\begin{proof}
The integral expansion in \(z-1\) defines the measure by Mahler duality.
We claim that
\[
 \mathcal F(z)+\mathcal F(-z)=0.
\]
The points \(z,-z\) are the two lifts over the same canonical degree-two quotient \(B\). The Serre--Tate pairing on étale and dual étale bases gives the relation \(z_{\rm target}^{e}=z_{\rm source}^c\) with the étale and connected factors \(e,c\) of the isogeny. 

Trace along the connected quotient morphism vanishes by expansion at Tate cusps (all exponents odd). More explicitly one does the trace on the tame ordinary curve, transporting tame data by the quotient. This is a finite-flat square correspondence: the quotient map reduces to Frobenius on the identified special tame curves, is locally a square map in deformation parameters, and the Tate map on cusp charts. Finiteness likewise lifts from the special Frobenius on the adic formal curves. Its trace identity on weight-zero functions follows by the same \(\mu\)-cusp detection (the two Tate expansions with opposite root arguments, and source cyclic choices still \(\mu\)-type). The left side is the transform of the measure whose multiplier is
\(1+(-1)^x\).  Uniqueness of the transform and torsion-freeness of the
coefficients show that the restriction of \(\mu\) to \(2\mathbb Z_2\)
is zero.  This proves the support assertion.
\end{proof}

\begin{lemma}[The exact logarithmic primitive]
\label{split-measures:logarithmic-primitive}
The modular logarithm \(G\) is analytic on each interior ordinary CM disk,
and on these disks
\[
 F_0=P_2(V_2)G.
 \tag{M1}\label{eq:M1}
\]
With the branch of the logarithm satisfying \(\log 2=0\), the Eisenstein
primitive is
\[
 F_{{\rm Eis},0}
   =-\frac1{48}\log\left(
      \frac{\Delta^2V_2^2\Delta}{(V_2\Delta)^3}\right).
\]
\end{lemma}
\begin{proof}
The log \(G\) used on such interior disks about CM reductions is analytic on the disk, with the prescribed group-log values at algebraic points. One justification, also for abelian projections, is by Néron extension on smooth charts. Near an interior ordinary tame point use the algebraic relative subgroup scheme of the fine elliptic moduli, i.e. the finite-flat-subgroup Hilbert moduli in the universal \(2^b\)-torsion, for subgroups of order \(2^b\) if the level requires it. It is of finite presentation and smooth over the trait at the canonical-subgroup special point: in Artin deformations there the subgroup must map trivially to the étale quotient (by étaleness and its special connectedness), so the unique lift over any ordinary elliptic deformation is the multiplicative subgroup. Thus the chart has the same smooth deformation germ as the tame chart. 

On its smooth open neighbourhood our characteristic-zero map to the abelian variety extends by the Néron property (in the generic fiber use the map on the cyclic locus; other subgroup types, if present, are open and closed and can use a constant). One may enlarge the local base first. Thus the disk maps into a single Néron tube, where log differences from the center are analytic by the formal group. On smaller closed disks one may first multiply into the small log range, giving the same analytic group log by division. This suffices for all level choices and their quotient disks here.

To prove \eqref{eq:M1}, compare differentials.  They agree by the parametrization differential, quotient pullbacks and the expansion normalization above (\(V_2 G\) differentiates with the extra Tate factor 2, \(V_2^2G\) with factor 4; these differential identities also follow by expansion on the ordinary curve). The possible constant is zero by summing over \(z,-z\). Indeed the two sources can be seen as the two quotients of \(B\) by noncanonical order-two subgroups, with inherited cyclic tame data. The \(T_2\) or \(U_2\) action gives
\[
                      G(z)+G(-z)=a_2 G(B)-\epsilon_2 G(V_2 B).
\]
For level at 2 these two noncanonical kernels are exactly the \(U_2\) quotients, retaining canonical \(2\)-level on the result. This Hecke normalization on divisors pulls back the usual differential operator on forms; cuspidal terms are torsion. Thus \(P_2(V_2)G\) also has trace zero in the pair.

For the Eisenstein assertion, roots of unity are killed by the logarithm
and \(\log 2=0\).  The displayed weight-zero discriminant ratio is nonvanishing on the interior disks, with analytic log (the normalized variation of an analytic unit on an open disk from its center has values in principal units, e.g. by its Newton polygon). The differentials agree by the log derivative of \(\Delta\). The norm of the ratio in the source pair is one by the Tate identity \(\Delta(q)\Delta(-q)=-\Delta(q^2)^3/\Delta(q^4)\); use the same correspondence expansion (the Tate orders of the ratio cancel at canonical cusps). This fixes the remaining disk constant and proves the Eisenstein identity.
\end{proof}

\subsection{Ring-class orbits and interpolation}

We next assemble the disk measures along a ring-class orbit. Use \(s=1\) or \(s=r_i\) as odd conductor here; the construction works also at odd products of auxiliary good primes coprime to the split levels, with conductor-\(s\) ring class weights denoted by \(\psi_s\) (here trivial or \(\Psi_{u,i}\)). In this description \(K\) can more generally have units just signs, split at the support \(2,S^\circ\), and \(\Gamma,\Psi_t\) defined as above (inertia at the 2's of finite index by the class sequence). Fix a CM point of order \(\mathcal O_s\), with the chosen split cyclic tame
orientations, and its translates indexed by
\(\mathfrak a\in\operatorname{Pic}(\mathcal O_s)\).

\begin{lemma}[CM disk coordinates]
\label{split-measures:cm-coordinates}
The conductor-\(s\) CM points are canonical lifts.  In their Serre--Tate
disks, the points \(z=\zeta_{2^m}^j\), for \(j\) odd, have endomorphism
order \(\mathcal O_{2^ms}\) and are defined over \(H_{2^ms}\) with the
specified cyclic level.  The kernel of the order-class map to conductor
\(s\) acts on them through the unit ratio
\((\mathbb Z/2^m\mathbb Z)^\times\), simply transitively on the primitive
roots in each disk.  Its inverse-limit map to \(\Gamma\) is a fixed map
\(\eta:\mathbb Z_2^\times\to\Gamma\), independent of \(s\).
\end{lemma}
\begin{proof}
The conductor-\(s\) points are canonical lifts because reduction is ordinary at the split prime, the endomorphisms split the \(2\)-divisible directions, and the Serre--Tate parameter must be one by the endomorphism lifting relation. Consequently the geometric special end ring is exactly \(\mathcal O_s\) (all such endomorphisms lift at the trivial parameter). Make the disk evaluations over the completed unramified base, enlarging constants if needed for the form coefficients. Choose Serre--Tate generators in both directions by polarization. For each \(\mathfrak a\) use a prime-to-\(2s\) and prime-to-level ideal action as representative, viewed also at higher 2-conductors. Write \(\delta_{\mathfrak a}\in\mathbb Z_2^\times\) for its Serre--Tate exponent from the base (connected divided by étale factor); take the identity at the trivial label.

Points at \(z=\zeta_{2^m}^{\,j}\), \(j\) odd for a primitive \(\zeta_{2^m}\), have endomorphisms \(\mathcal O_{2^m s}\): among endomorphisms of the reduction the two local actions must agree modulo \(2^m\). Canonical cyclic \(2\)-level is preserved by the order, as are the oriented tame groups. Thus by the main theorem of CM these \(X_0\)-points are defined over \(H_{2^m s}\): the units of the completed endomorphism order stabilize the cyclic level. 

Prime-to-level ideal actions transport the groups (including the canonical group at 2) and the disk parameters by the isogeny exponent. In particular the kernel of lowering the order class map to conductor \(s\), the unit ratios \((\mathbb Z/2^m)^\times\), acts simply transitively on the primitive roots in a disk. 

For an ideal principal below, extending to \((\gamma)\), the exponent is the local ratio \(\gamma_w/\gamma_{\bar w}\) up to reversing conventions (action by the corresponding endomorphism on the underlying conductor-\(s\) reduction). Thus in the inverse limit the exponent on the order-class kernel maps to \(\Gamma\) by one fixed \(\eta:\mathbb Z_2^\times\to\Gamma\), independent of \(s\) by the order class sequence. Here isogeny actions in all labels are identified by their actual Artin actions on the CM orbit, inverting ideals consistently if required by reciprocity. The cyclic \(2\)-level throughout is the canonical connected subgroup;
there is no independent variation of a \(2\)-level structure along the
orbit.
\end{proof}

Fix henceforth the Serre--Tate generators and compatible ideal
representatives used in Lemma~\ref{split-measures:cm-coordinates}.  Write
\(\delta_{\mathfrak a}\in\mathbb Z_2^\times\) for their connected-to-étale
isogeny exponents, with \(\delta_1=1\).

Write \([y]\) for the group-like \(t\)-power series of \(y\in\Gamma\),
using additive notation within \(\Gamma\). Let \(\mu_{\mathfrak a}\) be
the measure of \(F\) on the corresponding disk.  The chosen ideal
representative acts compatibly at all higher dyadic conductors; let
\(t_{\mathfrak a}\in\Gamma\) be the image of this Artin action in
\(\Gamma\).  With these fixed representatives, form
\[
 b_s(t,u)=\sum_{\mathfrak a}
       [\,t_{\mathfrak a}-\eta(\delta_{\mathfrak a})\,]\,\psi_s(\mathfrak a)
                         \int_{\mathbb Z_2^\times}[-\eta(x)]\,d\mu_{\mathfrak a}(x)
\]
(without \(u\) for \(s=1\)). This is bounded-integral, with conductor-\(s\) variables in finite cyclic group rings. Use it twice, with all odd orientations opposite but same base points and parameter choices on the underlying elliptic disks; write \(B_s=b_s^+ b_s^-\), \(B=B_1\), \(B_r=\lim B_{r_i}\) over the constant enlargement using \(\mathcal V^*\) in the integral series conventions.

\begin{proposition}[Character and central interpolation]
\label{split-measures:interpolation}
Let \(\theta\) be a sufficiently high finite character of \(\Gamma\), and
let \(\theta_l=\theta\circ\eta\) have primitive conductor \(2^m\).
Let \(X\) be the point with \(z=\zeta_{2^m}\) in the base disk, at the
original primitive level, and put \(\chi=\theta\psi_s\).  Then
\[
 \left(\sum_{j\ {\rm odd}\bmod 2^m}\theta_l(j)\zeta_{2^m}^{j}\right) b_s(\theta,u)
     =
 \left(\sum_{\rho\in\operatorname{Pic}(\mathcal O_{2^m s})}
                        \chi(\rho)\,G(X^\rho)\right)
                         \prod_{\ell\in S^\circ}P_\ell(\chi(\sigma_\ell)).
                                                                    \tag{M2}\label{eq:M2}
\]
Here \(\sigma_\ell\) is the inverse class translation for the oriented
\(V_\ell\), namely one of the two Frobenius classes.  At the trivial
character in \(t\), the corresponding identity is the sum over
\(\operatorname{Pic}(\mathcal O_s)\) of the disk-center logarithms, with
its odd Euler factors and the additional factor
\(P_2(\psi_s(\sigma_2))\).  For the elliptic curve \(E\), when
\(S^\circ\) is precisely the set of odd primes dividing \(N\), this gives
\[
 \begin{gathered}
 B(0)=\pm c_E^{-2}\log_{\omega_E}(P)^2
       \prod_{\ell\mid2N}P_\ell(1)^2,\\
 P=\sum_{\rho\in\operatorname{Pic}(\mathcal O_K)}
       \phi(X_1^\rho)\in E(K).
 \end{gathered}
 \tag{M3}\label{eq:M3}
\]
where \(X_1\) is a split Heegner point of conductor one.
\end{proposition}
\begin{proof}
The class-group exact sequence makes \(\theta\) factor through conductor
\(2^m\).  Fourier summation over odd \(j\) gives
\eqref{eq:M2}: by Lemma~\ref{split-measures:unit-support}, its nonzero
Gauss sum multiplies the inverse-character moment of each disk measure. The corrected weight within disk \(\mathfrak a\) at \(\zeta_{2^m}^{j}\) is \(\theta(t_{\mathfrak a})\theta_l(j/\delta_{\mathfrak a})\psi_s(\mathfrak a)\), precisely the full character value by the parameter description. Each tame \(V_\ell\) translates the orbit at raised level with the requisite orientations retained after lowering. In (M1), \(V_2,V_2^2\) disappear by primitivity (same quotient in the pair). 

The two primitive sums of \(G\) at opposite odd orientations differ only by a translation character unit and a sign: the Atkin--Lehner operation at the full odd part of the primitive level has this eigenlaw up to killed cusps, reverses all original odd choices and translates the orbit by the indicated odd ideals, preserving canonical level at 2.

At \(t=0\) there is instead the disk-center version: sum once over \(\operatorname{Pic}(\mathcal O_s)\) with weights \(\psi_s\); the same Euler translations occur at the odd primes, and an additional \(P_2(\psi_s(\sigma_2))\) at the inverse connected-isogeny translation by the canonical centers. This follows by augmentation and (M1). Taking \(s=1\), the two orientation sums differ by a sign and a
translation of the same Hilbert-class trace, giving \eqref{eq:M3}.
The factor \(c_E^{-2}\) is retained exactly; no assertion that \(c_E\) is
a 2-adic unit is used.
\end{proof}

\subsection{Trace compatibility and residual primitivity}

For each conductor \(r_i\), choose the base disk by a descending
order-\(r_i\) isogeny from the conductor-one base point, using a
non-eigenline.  Transport the Serre--Tate generators by polarization so
that the exponent is \(r_i^{\pm1}\).  In this subsection \(B_r\) is
formed using these compatible base disks.

\begin{lemma}[The split trace and residual order]
\label{split-measures:trace-residual}
Return to the geometric setting of
Proposition~\ref{split-measures:reducible-input}, without imposing its
Selmer-corank hypothesis, and take \(S^\circ\) to be exactly the odd primes dividing \(N\).  For the
auxiliary sequence of Lemma~\ref{split-measures:auxiliary-primes},
\[
 B_r(t,0)=(a_*-2)^2B(t),
 \tag{M4}\label{eq:M4}
\]
and
\[
 \operatorname{ord}_t\overline B
    =4\sum_{\substack{q\mid N\\q\ne2}}g_q.
 \tag{M5}\label{eq:M5}
\]
\end{lemma}
\begin{proof}
Return to the scalar residual test setting. For conductor \(r_i\) use a base disk chosen by an order-\(r_i\) descending isogeny from the conductor-one base (non-eigenline), transporting generators so the Serre--Tate exponent is \(r_i^{\pm1}\) by polarization. Thus for fixed \(m\) above, the evaluation basepoint is eventually also the corresponding \(r_i\)-descendant (the exponent is one modulo \(2^m\)). Trace down that conductor at \(u=0\) in (M2) gives multiplier \(a_{r_i}-2\) eventually, by the good-prime Hecke relation subtracting the two split prime-ideal translations, whose \(\theta\)-values become exactly 1. The two orientations give its square.  Applying the high-character
identity test to infinitely many such \(\theta\) proves \eqref{eq:M4}.

Finally compute modulo 2 at \(s=1\) by the Eisenstein congruence on the disks. The odd depletions on each measure contribute \(1-(q+1)[\pm b_q]/q+[\pm2b_q]/q\) at \(q\), by translation as above (test at high characters by the same Fourier summation), of residual order \(2g_q\). 

The measure using just \(F_{{\rm Eis},0}\) has unit value at the center, which is its unweighted CM trace. Indeed the product of \(V_2\Delta/\Delta\) over the conductor-one class orbit is \((\alpha/2^h)^{12}=\bar\alpha^{-12}\), viewing \(\alpha\) at the chosen embedding with \(w\) connected. The connected isogeny permutes the classes, normalizing differential pullback by 2 in \(V_2\); the ordinary isogeny pullback factors for independently chosen differentials have product the principal-generator factor around the cycles, altogether \(\pm\alpha\). Similarly the product of \(V_2^2\)-ratios is \(\bar\alpha^{-24}\). Thus the log-ratio trace with factor \(-1/48\) is \(-\log(\bar\alpha)/4\), odd by the dyadic basis congruence above. Thus the Eisenstein measure before the odd depletions is a unit at the center.  Each
odd depletion contributes order \(2g_q\) to each of the two orientations,
which proves \eqref{eq:M5}.
\end{proof}

The matching residual orders \eqref{eq:A1} and \eqref{eq:M5} will identify
an integral divisibility quotient as a unit.  The transverse test
\eqref{eq:A3} permits the horizontal divisibilities needed for that step
to be lifted back to the integral two-variable ring.
